% Source: https://artofproblemsolving.com/community/c3h2461075_was_ramanujans_estimate_for_pi_really_an_estimate
% Stable identity: aops-topic-2461075-q01
\textbf{(a)} $9^2+\dfrac{19^2}{22}=81+\dfrac{361}{22}=\dfrac{1782+361}{22}=\boxed{\dfrac{2143}{22}}$, in lowest terms since $2143$ is odd and not divisible by $11$ ($2143=11\cdot194+9$). As a decimal, $\dfrac{2143}{22}=97.40909090\ldots$, with $09$ repeating.

\textbf{(b)} Since $\pi>3.1415926535$,
\[\pi^2>3.1415926535^2=9.869604400525\ldots>9.8696044,\]
and therefore
\[\pi^4>9.8696044^2=97.40909101249936>97.4090909090\ldots=\dfrac{2143}{22}=r^4.\]
As both are positive, $r<\pi$. $\blacksquare$

\textbf{(c)} Similarly, $\pi^2<3.1415926536^2=9.869604401153\ldots<9.8696044012$, so
\[\pi^4<9.8696044012^2<97.4090910362,\]
and
\[0<\pi^4-r^4<97.4090910362-97.4090909090<1.3\cdot10^{-7}.\]
Now factor $\pi^4-r^4=(\pi-r)(\pi+r)\left(\pi^2+r^2\right)$. Both $\pi$ and $r$ exceed $3$ (because $r^4=97.4\ldots>81$), so $(\pi+r)\left(\pi^2+r^2\right)>6\cdot18=108$, and
\[0<\pi-r<\dfrac{1.3\cdot10^{-7}}{108}<1.3\cdot10^{-9}.\]
Hence
\[0<1-\dfrac r\pi=\dfrac{\pi-r}{\pi}<\dfrac{1.3\cdot10^{-9}}{3}<5\cdot10^{-10},\]
so $0.9999999995<\dfrac r\pi<1$. Rounded to $9$ significant digits this is $1.00000000$, which is why the calculator shows $1$. In fact $r=3.14159265258\ldots$ while $\pi=3.14159265358\ldots$: the approximation is correct to $8$ decimal places, but it is not exact.
